\documentclass[aspectratio=169]{beamer}

\usetheme{Madrid}
\usecolortheme{seahorse}

\usepackage[utf8]{inputenc}
\usepackage[T1]{fontenc}
\usepackage{lmodern}
\usepackage{booktabs}
\usepackage{graphicx}
\usepackage{array}
\usepackage{xcolor}

\title[Tear Desiccate Classification]{Analysis of Human Tear Desiccate\\ based on Data Analysis and Machine Learning}
\subtitle{HackKošice 2026 -- Pavol Jozef Šafárik University in Košice track}
\author{Luka Divjakovi\'c \and Veselin Roganovi\'c}
\institute{HackKošice}
\date{\today}

\begin{document}

\begin{frame}
    \titlepage
\end{frame}

\begin{frame}{Outline}
    \tableofcontents
\end{frame}

% =====================================================================
\section{Data Description and Preprocessing}
% =====================================================================

\begin{frame}{Dataset overview}
    \begin{itemize}
        \item Provided dataset: \textbf{5 folders}, one per class (chronic condition / healthy).
        \item Each folder contains \textbf{images} (\texttt{.bmp}) and \textbf{raw data} files.
        \item Dataset is \textbf{relatively small} -- we searched online for additional
              public data without success, but studied several scientific papers on the topic.
        \item After exploring both modalities, we decided to use \textbf{only the \texttt{.bmp} images}.
        \item Pipeline: we store each image's \textbf{file path} and the \textbf{class label}
              (from its parent folder) in a DataFrame; the actual image is loaded on demand
              during training / prediction.
    \end{itemize}
\end{frame}

\begin{frame}{Patterns observed per class}
    \small
    \begin{itemize}
        \item \textbf{Diabetes}: dry--leaf-like branching, many white dots, flake-like structures,
              sharp-angled straight lines.
        \item \textbf{Glaucoma}: unusual dots, droplets and spiky objects; mostly smooth curved lines,
              fewer sharp angles.
        \item \textbf{Sclerosis}: similar to diabetes but dots are more spaced; some regions are entirely white,
              others look flaky but are in fact clusters of dots spread everywhere.
        \item \textbf{Dry eye}: even more similar to diabetes; extremely regular dot groupings forming
              cypress-like shapes; very clean, well-defined clusters.
        \item \textbf{Healthy}: similar to glaucoma; more curved lines than dot clusters, occasional bright
              regions, cypress-like line patterns without visibly separated dots.
    \end{itemize}
\end{frame}

\begin{frame}{Data quality issues}
    \begin{itemize}
        \item A portion of images are \textbf{zoomed in / at different scales} -- may mislead the model.
        \item Presence of \textbf{duplicates and near-duplicates}.
        \item Some images are visibly \textbf{damaged / low quality}.
        \item Every image contains a \textbf{white frame and scale/metadata overlay}.
    \end{itemize}
    \vspace{0.5em}
    $\Rightarrow$ We removed a subset of clear outliers so the model can \textbf{generalize better}.
\end{frame}

\begin{frame}{Preprocessing pipeline}
    \begin{enumerate}
        \item \textbf{ROI crop}: remove the white frame and the scale/metadata overlay.
        \item \textbf{Grayscale conversion}: color carries little signal, and our algorithms
              performed better on grayscale.
        \item \textbf{CLAHE contrast normalization}: enhances local contrast and makes
              fine-grained tear-film patterns more discriminative.
        \item \textbf{Data augmentation -- rotations}: class patterns are not direction-specific,
              so rotating each image produces valid extra training samples and enlarges the dataset.
        \item \textbf{Outlier filtering}: drop visibly damaged / duplicate images before training.
    \end{enumerate}
\end{frame}

% =====================================================================
\section{Algorithms and Methods}
% =====================================================================

\begin{frame}{Evaluation protocol}
    \begin{itemize}
        \item \textbf{Train / test split}: 80\,/\,20, stratified by class.
        \item \textbf{Cross-validation}: 5-fold CV on the \textbf{train} portion for model selection
              and hyper-parameter tuning.
        \item \textbf{Test set}: used only \textbf{once}, after we are happy with the model on CV,
              to obtain the final score.
        \item \textbf{Metric}: weighted F1-score (as required by the challenge).
    \end{itemize}
\end{frame}

\begin{frame}{Classical approach I -- HOG + SVM / KNN}
    \begin{itemize}
        \item Images clearly contain \textbf{oriented gradient patterns} (lines, dendritic
              shapes, dot clusters).
        \item \textbf{Histogram of Oriented Gradients (HOG)} captures local edge
              directions in cells and pools them into a feature vector.
        \item The HOG descriptor is then fed into classical classifiers:
              \begin{itemize}
                  \item \textbf{SVM} with RBF/linear kernel,
                  \item \textbf{KNN} as a simple non-parametric baseline.
              \end{itemize}
        \item Motivation: interpretable, fast, and a natural fit for the type of
              texture/gradient patterns visible in the tear images.
    \end{itemize}
\end{frame}

\begin{frame}{Classical approach II -- KMeans ($K{=}3$) + SVM / KNN}
    \begin{itemize}
        \item Visually, most images contain roughly \textbf{3 distinct regions}:
              \begin{itemize}
                  \item background,
                  \item dots / main figure,
                  \item bright white dots.
              \end{itemize}
        \item We use \textbf{KMeans with $K{=}3$} on pixel intensities to segment
              the image into these regions.
        \item Statistics of the resulting clusters (area ratios, intensities,
              spatial descriptors) form a compact feature vector.
        \item This vector is classified by \textbf{SVM} or \textbf{KNN}.
    \end{itemize}
\end{frame}

\begin{frame}{Deep learning -- Transfer learning with pretrained CNNs}
    \begin{itemize}
        \item The dataset is too small to train a CNN from scratch.
        \item We use \textbf{ImageNet-pretrained} backbones, replace the classification head,
              and \textbf{unfreeze the last few blocks} for fine-tuning.
        \item Two-stage training: \textbf{head only}, then \textbf{partial fine-tuning}.
        \item Models tried: \texttt{ResNet18}, \texttt{ResNet50}, \texttt{EfficientNet-B0},
              \texttt{DenseNet121}, \texttt{MobileNetV3-Large}, \texttt{ViT-B/16}.
        \item The four listed first gave the best results.
        \item Each image is CLAHE-normalized, replicated into 3 channels,
              resized to $224{\times}224$ and standardized with ImageNet mean/std.
    \end{itemize}
\end{frame}

% =====================================================================
\section{Analytical Objectives}
% =====================================================================

\begin{frame}{Analytical objectives}
    Our methods address several analytical tasks on the tear-desiccate images:
    \begin{itemize}
        \item \textbf{Pattern recognition}: identifying class-specific morphological
              patterns (branching dendrites, dot clusters, cypress-like structures).
        \item \textbf{Recurring structure identification}: KMeans segmentation
              exposes recurring region types (background, figure, bright dots).
        \item \textbf{Similarity detection}: duplicate / near-duplicate detection
              during dataset cleaning.
        \item \textbf{Data comparison}: inter-class comparison of gradient
              statistics (HOG) and region statistics (KMeans).
        \item \textbf{Feature learning}: CNN backbones learn high-level
              representations that encode correlations between fine-grained
              texture patterns and the underlying chronic condition.
    \end{itemize}
\end{frame}

% =====================================================================
\section{Expected Outcome and Results}
% =====================================================================

\begin{frame}{Task and expected outcome}
    \begin{itemize}
        \item \textbf{Goal}: a classifier that, given a tear-desiccate image,
              estimates the probability that a patient has a specific
              \textbf{chronic disease}.
        \item \textbf{Evaluation}: a hidden test set; ranking by
              \textbf{weighted F1-score}.
        \item \textbf{Winner}: the team with the highest weighted F1 wins the track.
    \end{itemize}
\end{frame}

\begin{frame}{Results}
    \centering
    \begin{tabular}{clcc}
        \toprule
        \textbf{\#} & \textbf{Model} & \textbf{CV F1 (weighted)} & \textbf{Test F1 (weighted)} \\
        \midrule
        1 & \textbf{ResNet50} \textit{(winner)} & \textbf{0.8220} & \textbf{0.7015} \\
        2 & ResNet18                             & 0.7973 & 0.5854 \\
        3 & DenseNet121                          & 0.7217 & 0.5834 \\
        4 & EfficientNet-B0                      & 0.7096 & 0.6849 \\
        5 & KMeans ($K{=}3$) + SVM               & 0.6974 & 0.5582 \\
        6 & HOG + SVM                            & 0.4916 & 0.5155 \\
        \bottomrule
    \end{tabular}

    \vspace{1em}
    \small
    Deep transfer learning clearly outperforms classical pipelines.
    \textbf{ResNet50} is our final submitted model.
\end{frame}

\begin{frame}{Conclusion}
    \begin{itemize}
        \item Careful preprocessing (ROI crop, grayscale, CLAHE, rotation augmentation,
              outlier removal) was crucial given the small dataset.
        \item Classical HOG/KMeans pipelines provide interpretable baselines
              but are limited by the handcrafted features.
        \item Transfer learning with pretrained CNNs dominates both CV and test
              performance; \textbf{ResNet50} achieves the best weighted F1.
        \item Future work: more data, ensembling across backbones, self-supervised
              pretraining on raw tear images.
    \end{itemize}
\end{frame}

\begin{frame}
    \vfill
    \centering
    {\Huge Thank you!}\\[1em]
    {\large Questions?}\\[2em]
    \small Luka Divjakovi\'c \quad\&\quad Veselin Roganovi\'c \\
    HackKošice -- Pavol Jozef Šafárik University in Košice track
    \vfill
\end{frame}

\end{document}
